\section{Results III: where the labels should go}
\label{sec:where}

\Cref{tab:where} puts every way of using 5,000 labels that we tested on one scale: the test accuracy of the
final classifier, with only those labels.

\begin{table}[t]
\centering
\small
\caption{Where the labels enter (CIFAR-10, 5,000 labels, ResNet-18). Test accuracy (\%) of the final
classifier, mean $\pm$ standard deviation over seeds. Pretraining is 19,400 steps where it applies.
FixMatch and MixMatch from scratch process as many views as a 9,700-step contrastive run, or three times as
many where stated; FixMatch after pretraining adds the same amount to the 19,400 pretraining steps.}
\label{tab:where}
\begin{adjustbox}{max width=\linewidth}
% generated by scripts/make_assets.py -- do not edit
\begin{tabular}{llcc}
\toprule
pipeline & where the labels enter & seeds & test accuracy \\
\midrule
\multicolumn{4}{l}{\emph{labels inside the contrastive loss from the first step}} \\
\quad \supcon, then FT-5k & inside the loss & 3 & 82.2{\scriptsize$\,\pm\,$0.2} \\
\quad \supcon, then FT-5k (oracle recipe) & inside the loss & -- & -- \\
\quad \union, then FT-5k & inside the loss & 3 & 81.7{\scriptsize$\,\pm\,$0.3} \\
\quad \union, then FT-5k (oracle recipe) & inside the loss & -- & -- \\
\multicolumn{4}{l}{\emph{labels after label-free pretraining}} \\
\quad \simclr, then LP-5k & after & 5 & 85.8{\scriptsize$\,\pm\,$0.2} \\
\quad \simclr, then FT-5k & after & 3 & 88.6{\scriptsize$\,\pm\,$0.4} \\
\quad \simclr, then FT-5k (selected recipe) & after & 3 & 89.4{\scriptsize$\,\pm\,$0.1} \\
\quad \simclr, then FT-5k (oracle recipe) & after & -- & -- \\
\quad BYOL, then FT-5k & after & -- & -- \\
\quad \simclr, then \supcon on the 5k (485 steps), then LP-5k & after, contrastive & 3 & 87.7{\scriptsize$\,\pm\,$0.3} \\
\quad \simclr, then \supcon on the 5k (485 steps), then FT-5k & after, contrastive & -- & -- \\
\quad \simclr, then FT-5k, then self-distillation & after & -- & -- \\
\multicolumn{4}{l}{\emph{labels in a semi-supervised objective}} \\
\quad FixMatch from scratch & pseudo-labels & 1 & 90.8 \\
\quad FixMatch from scratch, 3$\times$ budget & pseudo-labels & -- & -- \\
\quad MixMatch from scratch & pseudo-labels & 1 & 92.4 \\
\quad \simclr, then FixMatch & after, pseudo-labels & -- & -- \\
\quad BYOL, then FixMatch & after, pseudo-labels & -- & -- \\
\quad \supcon, then FixMatch & both & -- & -- \\
\quad \union, then FixMatch & both & -- & -- \\
\quad \simclr, then \supcon (2,425 steps), then FixMatch & both & -- & -- \\
\multicolumn{4}{l}{\emph{references}} \\
\quad cross-entropy on the 5k, 9,700 steps & supervised only & 3 & 82.6{\scriptsize$\,\pm\,$0.3} \\
\quad cross-entropy on the 5k, 19,400 steps & supervised only & 3 & 83.5{\scriptsize$\,\pm\,$0.5} \\
\quad cross-entropy on all 50k labels & all labels & 1 & 94.8 \\
\bottomrule
\end{tabular}

\end{adjustbox}
\end{table}

\paragraph{Labels inside the loss from the first step}
\supcon followed by fine-tuning reaches \R{supcon_lab_f0.1_e200:ft:pm}\%, and \union
\R{union_f0.1_e200:ft:pm}\%. Plain cross-entropy training on the same 5,000 images for the same 19,400
steps reaches \R{ce_f0.1_e200:clf:pm}\%. With the default fine-tuning recipe, contrastive pretraining with
the labels is therefore no better than not pretraining at all
(\R{cmp:supconft-ce:e200:sdiff}\pt, \R{cmp:supconft-ce:e200:ci}).

\paragraph{Labels after label-free pretraining}
\simclr followed by fine-tuning reaches \R{simclr_all_e200:ft:pm}\%, \R{cmp:simclrft-ce:e200:diff}\pt above
cross-entropy on the 5,000 labels. Two further arms keep the labels until after pretraining.

\emph{\simclr, then \supcon.} We continue the \simclr encoder with \supcon on the 5,000 labeled images, with
a small learning rate (0.02) and for a short time. Five passes of the full-data schedule (485 steps) raise
\lpfive from \R{simclr_all_e200:lp5}\% to \R{twostage_simclr200_supcon5:lp5:pm}\%
(\R{cmp:twostage-simclr:lp5:e200:ci}) and \lpall from \R{simclr_all_e200:lp50}\% to
\R{twostage_simclr200_supcon5:lp50}\%. Longer second stages are worse (\R{twostage_simclr200_supcon25:lp5}\%
\lpfive after 2,425 steps), and with the first-stage learning rate of 0.25 the benefit of pretraining is
largely lost. The labels thus help inside a contrastive loss when they arrive late and briefly; the same
loss applied from the first step ends \R{d:twostage-supcon:lp5}\pt lower. \pending{Fine-tuning after the
short \supcon stage gives \R{twostage_simclr200_supcon5:ft:pm}\%.}

\emph{\simclr, then fine-tuning, then self-distillation.} In the style of SimCLR~v2 \citep{chen2020simclrv2},
the fine-tuned network labels all 50,000 training images with class probabilities and a student, started
from the same \simclr encoder, is trained to match them. \pending{This reaches
\R{simclr_all_e200:dist:pm}\%.}

\paragraph{Labels in a semi-supervised objective}
FixMatch and MixMatch trained from scratch reach \R{fixmatch_f0.1:clf:pm}\% and
\R{mixmatch_f0.1:clf:pm}\% with the compute of a 9,700-step contrastive run. They use the labels from the
first step and still beat every pretrain-then-fine-tune pipeline in the table, including \simclr with
twice the pretraining budget. \pending{Starting FixMatch from the \simclr encoder gives
\R{simclr_all_e200:fm:pm}\%, the best result of our study on CIFAR-10 with 5,000 labels; FixMatch from
scratch with three times the budget, a comparable total compute, gives \R{fixmatch_b3_f0.1:clf:pm}\%.}

\paragraph{Summary}
The answer to the question in the title is ``after''. Within the pipelines that end with supervised
fine-tuning, the labels are better kept out of the contrastive loss until a label-free encoder exists.
Overall, the best place for scarce labels in our experiments is not a contrastive loss at all but a
pseudo-labeling objective, alone or after label-free pretraining. This is consistent with the semi-supervised
literature \citep{sohn2020fixmatch,li2021comatch,zheng2022simmatch,wang2022usb}, and it bounds the claim of
this paper: we compare ways of using labels around contrastive pretraining, and we do not claim that
contrastive pretraining is the best use of a label budget.
